% DermaJEPA preprint — pass-2 strict pre-publication review
% Three-agent parallel deep audit: zero-context claim re-derivation,
% adversarial / counter-narrative, mathematical formalism rigor.
% Date: 2026-04-28 14:36 UTC. Reviewer: paper-review skill.

\documentclass[11pt,a4paper]{article}
\usepackage[margin=2.5cm]{geometry}
\usepackage{enumitem}
\setlist{noitemsep, topsep=2pt}
\usepackage{hyperref}
\usepackage{xcolor}

\title{Pre-publication review (pass 2) of \emph{Frozen-backbone JEPA-style
       probes on HAM10000} (commit \texttt{365b641})}
\author{Three-agent deep audit (zero-context $\parallel$ adversarial $\parallel$ formalism)}
\date{2026-04-28}

\begin{document}
\maketitle

\section*{Headline}

Twelve headline numbers re-derived independently from the raw evidence
without paper-side context: \textbf{all verified}. The pass surfaces
\textbf{zero blockers}, \textbf{five major}, and \textbf{seven minor}
issues. The major issues are reviewer-anticipation defences and
exposition gaps, not factual errors. The math and the implementation
agree on every checkable point. Recommendation: \textbf{minor revision}
with the following inserted before submission.

\section{Major findings}

\subsection*{B1 (transparency) --- AUROC convention is sign-locked;
            symmetric-reading defence is missing}

\textbf{Location.} Abstract; \S5; Table~1; \texttt{baselines.py:44}
(\texttt{max(...,~key=auroc)}).\\
\textbf{Problem.} Under the paper's stated convention (label
\texttt{changing}~$=1$, score~$=$~cosine distance, no sign flip), raw
DermLIP cosine has $\auroc = 0.109$. A reviewer will note that
$1 - 0.109 = 0.891$ is the symmetric-AUROC equivalent and ask why
pixel-L2 ($0.580$) is the ``strongest baseline'' rather than
sign-corrected raw DermLIP cosine. Under that reading the JEPA scaffold
contributes only $+0.053$ on top of raw cosine, not $+0.363$.\\
\textbf{Why it matters.} The defence is real: under fixed train/test
convention the predictor must \emph{learn the correct sign}, which
cannot be flipped in deployment. But the paper never states this. A
TMLR / MIDL reviewer will treat the comparison as ill-posed without
the explicit clarification.\\
\textbf{Fix.} Add to \S3.5 (Evaluation protocol) a sentence stating
that $\auroc$ is computed directionally under the fixed convention; in
\S5 (or as a footnote on Table 1) report the symmetric-equivalent raw
embedding cosine $\auroc$ alongside the directional value, and a one-
sentence justification of the directional comparison: \emph{the JEPA
predictor's job includes learning the correct orientation; a baseline
that requires an a-priori sign flip is not a usable reference for our
probe.}

\subsection*{C-F2 (DE) --- predictor domain mismatch}

\textbf{Location.} \texttt{03-method.tex:21}.\\
\textbf{Problem.} $g_\theta : \mathbb{R}^d \to \mathbb{R}^d$, but the
input is always $z(x) \in \mathbb{S}^{d-1}$.\\
\textbf{Fix.} Replace with $g_\theta : \mathbb{S}^{d-1} \to \mathbb{R}^d$.
The codomain stays $\mathbb{R}^d$ (output is L2-normalised only at
scoring time, eq.~3).

\subsection*{C-F5 (DE) --- AUROC label convention never stated}

\textbf{Location.} \texttt{03-method.tex:50--52}.\\
\textbf{Problem.} The paper says ``a correctly oriented predictor
produces $s(x_c, x_t) < s(x_c', x_t')$ whenever $(x_c, x_t)$ is
stable'' but never states which class is positive in the AUROC
calculation. Verified via \texttt{contracts.py:44--45} and
\texttt{training.py:448}: positive class is \texttt{changing}~$=1$;
score is the raw cosine distance (no sign flip).\\
\textbf{Fix.} Insert after eq.~3: ``We adopt $y = 1$ for changing and
$y = 0$ for stable; $s$ is used directly (no sign flip) as the AUROC
score, so a correctly oriented predictor yields $\auroc > 0.5$.''

\subsection*{C-F8 (UC) --- ``no transform type shared'' is false read
            literally}

\textbf{Location.} \texttt{03-method.tex:71--93} and Appendix~C.\\
\textbf{Problem.} JPEG compression appears in all three families
($q \in \{45\text{--}70, 20\text{--}40, 80\text{--}95\}$). The
``operation-level disjointness'' claim is not literal disjointness of
operation \emph{types}; it is disjointness of operation--parameter-band
\emph{combinations}.\\
\textbf{Fix.} ``no operation--parameter-band combination is shared
between any two families; JPEG is the one operation that recurs across
all three, with disjoint quality bands $\{20\text{--}40,
45\text{--}70, 80\text{--}95\}$, chosen so the resulting compression
artefact is perceptually distinct per family.''

\subsection*{B2 / B3 (scope softening) --- ``any medical-image
            pretraining'' and ``two architectures''}

\textbf{Location.} Abstract, \S1, \S5, \S7.\\
\textbf{Problem.} (a) ``Rules out any medical-image pretraining'' is
stronger than EXP-008 supports: PMC-15M's dermatology fraction is
undocumented in the paper, so the partition really only rules out
``PMC-15M-style biomedical-figure pretraining at the dermatology
fraction PMC-15M happens to contain''. (b) ``Two architectures''
overstates a comparison that mixes a DINOv2-pretrained ViT-B/14 with
three CLIP-pretrained ViT-B/16 backbones; better phrasing is ``two
pretraining paradigms within the ViT-B class.''\\
\textbf{Fix.} Reword in abstract / \S1 / \S5 / \S7 conclusions to the
narrower scope.

\section{Minor findings}

\begin{enumerate}[leftmargin=2em]
\item \textbf{C-F1.} $\mathcal{X}$ undefined. Add ``\,$\mathcal{X}
      \subset [0, 1]^{3 \times H \times W}$ denotes RGB images at the
      backbone's native input resolution''.
\item \textbf{C-F3.} Train/eval normalisation asymmetry is intentional
      but unstated. Add: ``We do not L2-project predictions during
      training; the loss
      $\|g_\theta(z(x_c)) - z(x_t)\|_2^2 = \|g_\theta(z(x_c))\|_2^2 -
      2\langle g_\theta(z(x_c)), z(x_t)\rangle + 1$ jointly drives the
      prediction direction toward $z(x_t)$ and the prediction norm
      toward $1$. Appendix~E reports prediction-norm mean $\geq 0.97$
      in every run.''
\item \textbf{C-F4.} ``Soft regulariser'' equilibrium argument missing.
      Add: ``$\mathcal{R}$ is soft, not constraining: at equilibrium
      $\|W - I\|_F$ is set by the balance between data gradient and
      $\lambda$, not pinned at zero.''
\item \textbf{C-F6.} Bootstrap is not label-stratified; the
      stable/changing balance per resample drifts $\sim \pm 22$. Add a
      sentence to Appendix~D acknowledging this and stating the CI is
      marginal over the empirical class balance.
\item \textbf{C-F7.} ``Collapse'' minima not specified. In Appendix~E
      define: ``prediction-norm minimum $\equiv \min_{x \in
      \mathcal{D}_{\mathrm{test}}} \|g_\theta(z(x))\|_2$;
      dimension-variance minimum $\equiv \min_{j} \mathrm{Var}_{x \in
      \mathcal{D}_{\mathrm{test}}}(g_\theta(z(x))_j)$''.
\item \textbf{C-F10.} $\tilde g_\theta$ should be defined in a numbered
      display rather than inline so \S7 can reference it.
\item \textbf{(Agent-A observation, not a defect.)} The
      general-medical-to-dermoscopy step is more precisely $+0.6149$
      (rounded to $+0.62$). The abstract's $+0.62$ rounds in the
      conventional direction; consider quoting $+0.61$ in any future
      precision-sensitive context.
\end{enumerate}

\section{Findings B could not break}

\begin{itemize}[leftmargin=2em]
\item Predictor near-identity vs train $\auroc = 0.953$: weights from
      \texttt{jepa\_predictor.npz} show $\|W - I\|_F = 2.33$ for
      EXP-007 and $3.49$ for EXP-004 with bias norm $\approx 0.44$.
      With the soft regulariser and the empirical loss gradient, this
      departure is compatible with the data; not a paradox.
\item EXP-002 (train $\auroc 0.953$) vs EXP-004 (train $\auroc 0.900$):
      different proxies (matched-eval vs mixed-family training); not a
      contradiction.
\item Seed-std ratio (DermLIP $0.0029$ vs BiomedCLIP $0.0120$):
      $F = 16.72$, $\mathrm{df}(4,4)$, two-sided $p = 0.018$. The
      $4\times$ ratio is statistically significant at $\alpha = 0.05$
      even with $n = 5$.
\item Lesion-ID-aware leakage probe: existing probe in Appendix~H is
      honestly scoped; Agent~B did not test diagnosis-correlation or
      photographer-identity leakage but did not invalidate the existing
      probe either.
\end{itemize}

\section{Top-priority fixes (top 5)}

\begin{enumerate}[leftmargin=2em]
\item B1 --- explicit AUROC-convention statement plus symmetric-equivalent
      raw-cosine row in Table~1 caption.
\item C-F2 --- correct $g_\theta$ domain to $\mathbb{S}^{d-1}$.
\item C-F5 --- state the label convention ($y = 1$ for changing,
      $y = 0$ for stable; no sign flip).
\item C-F8 --- soften ``no transform type shared'' to
      ``no operation--parameter-band combination shared''.
\item B2 / B3 --- scope softening on ``any medical-image pretraining''
      and ``two architectures''.
\end{enumerate}

\end{document}
